% GENERATED FILE. Edit canonical lessons, then run npm run generate:books.
% canonical-source-hash: fnv1a64:774bbe4fcda72433
% canonical-lessons: ML-C417-cattam, ML-C417-palikkuka, ML-C417-anusarikkuka, ML-C417-niyamaparamaya, ML-C417-niyamaviruddhamaya

\chapter{Rules and the Law}
\label{ch:ml-rules-and-the-law}
\hlchaptermodality{417}

\begin{chapteropening}
\textbf{By the end of this chapter:} \emph{I can say a rule, to abide by, to obey, legal and illegal.}

Complete the last lesson of chapter 417: I can say a rule, to abide by, to obey, legal and illegal.
\end{chapteropening}

\section[\texorpdfstring{\ml{ചട്ടം} (caṭṭaṁ)}{ചട്ടം (caṭṭaṁ)}]{\ml{ചട്ടം} (caṭṭaṁ) — a regulation; a rule}

\label{lesson:ML-C417-cattam}



\begin{quote}
Before the new one: say the Malayalam for human rights, then the Malayalam for discrimination.
\end{quote}

\subsection*{You'll want to know: \ml{ചട്ടം}}
\textbf{\ml{ചട്ടം}} — \emph{caṭṭaṁ} — \textquotedblleft{}a regulation; a rule\textquotedblright{}.

\begin{quote}
\ml{ഓഫീസിലെ} \ml{ചട്ടം} \ml{അനുസരിച്ച്} \ml{ഫോൺ} \ml{ഉപയോഗിക്കാൻ} \ml{പാടില്ല}.
\end{quote}

\begin{quote}
\emph{ōphīsile caṭṭaṁ anusariccŭ phōṇ upayōgikkān pāṭilla.}
\end{quote}

\begin{quote}
Under the office rules, phones may not be used.
\end{quote}

Plural: \ml{ചട്ടങ്ങൾ}, \emph{caṭṭaṅṅaḷ}.

Plural: \ml{ചട്ടങ്ങൾ}, regulations.

\begin{cousinweb}
\textbf{\ml{ചട്ടം}} (\emph{caṭṭaṁ}) is native Dravidian: DEDR 2304 lists Malayalam \emph{caṭṭam} as \textquotedblleft{}a frame, a mould, a plan, a regulation, a rule, a law\textquotedblright{}. A rule is the frame things are made in.
\end{cousinweb}

\begin{grammarlens}[title={What you've built}]
One more word for giving reasons and opinions.
\end{grammarlens}

\subsection*{Your turn}
\begin{itemize}
\raggedright
  \item \emph{Say it:} \emph{caṭṭaṁ}
  \item \emph{Say it:} \emph{caṭṭaṁ}, once more
  \item \emph{Say it:} \emph{ōphīsile caṭṭaṁ anusariccŭ phōṇ upayōgikkān pāṭilla.} — Under the office rules, phones may not be used.
  \item \emph{Recall:} say the Malayalam for human rights, then the Malayalam for discrimination, then say \emph{caṭṭaṁ} again
\end{itemize}

\begin{tcolorbox}[breakable,colback=teal!4,colframe=teal!35!black,title={Before you move on}]
What does \ml{ചട്ടം} mean? (\textquotedblleft{}A regulation; a rule\textquotedblright{}.) Say it once more.
\end{tcolorbox}
\section[\texorpdfstring{\ml{പാലിക്കുക} (pālikkuka)}{പാലിക്കുക (pālikkuka)}]{\ml{പാലിക്കുക} (pālikkuka) — to abide by; to keep (a rule)}

\label{lesson:ML-C417-palikkuka}



\begin{quote}
Before the new one: say the Malayalam for discrimination, then the Malayalam for a regulation.
\end{quote}

\subsection*{You'll want to know: \ml{പാലിക്കുക}}
\textbf{\ml{പാലിക്കുക}} — \emph{pālikkuka} — \textquotedblleft{}to abide by; to keep (a rule)\textquotedblright{}.

\begin{quote}
\ml{എല്ലാവരും} \ml{ട്രാഫിക്} \ml{നിയമങ്ങൾ} \ml{പാലിക്കണം}.
\end{quote}

\begin{quote}
\emph{ellāvaruṁ ṭrāphikŭ niyamaṅṅaḷ pālikkaṇaṁ.}
\end{quote}

\begin{quote}
Everyone must keep the traffic rules.
\end{quote}

Past: \ml{പാലിച്ചു}, \emph{pāliccu}.

A promise is also \ml{പാലിക്കുക}: \ml{വാക്ക്} \ml{പാലിക്കുക}, to keep one's word.

\begin{cousinweb}
\textbf{\ml{പാലിക്കുക}} (\emph{pālikkuka}) is built from Sanskrit \emph{pālana}, \textquotedblleft{}protecting, keeping\textquotedblright{} (CDIAL 8128), with the Malayalam verb ending -\ml{ിക്കുക}. To keep a rule is to guard it.
\end{cousinweb}

\begin{grammarlens}[title={What you've built}]
One more word for giving reasons and opinions.
\end{grammarlens}

\subsection*{Your turn}
\begin{itemize}
\raggedright
  \item \emph{Say it:} \emph{pālikkuka}
  \item \emph{Say it:} \emph{pālikkuka}, once more
  \item \emph{Say it:} \emph{ellāvaruṁ ṭrāphikŭ niyamaṅṅaḷ pālikkaṇaṁ.} — Everyone must keep the traffic rules.
  \item \emph{Recall:} say the Malayalam for discrimination, then the Malayalam for a regulation, then say \emph{pālikkuka} again
\end{itemize}

\begin{tcolorbox}[breakable,colback=teal!4,colframe=teal!35!black,title={Before you move on}]
What does \ml{പാലിക്കുക} mean? (\textquotedblleft{}To abide by; to keep (a rule)\textquotedblright{}.) Say it once more.
\end{tcolorbox}
\section[\texorpdfstring{\ml{അനുസരിക്കുക} (anusarikkuka)}{അനുസരിക്കുക (anusarikkuka)}]{\ml{അനുസരിക്കുക} (anusarikkuka) — to obey}

\label{lesson:ML-C417-anusarikkuka}



\begin{quote}
Before the new one: say the Malayalam for a regulation, then the Malayalam for to abide by.
\end{quote}

\subsection*{You'll want to know: \ml{അനുസരിക്കുക}}
\textbf{\ml{അനുസരിക്കുക}} — \emph{anusarikkuka} — \textquotedblleft{}to obey\textquotedblright{}.

\begin{quote}
\ml{കുട്ടികൾ} \ml{അധ്യാപകനെ} \ml{അനുസരിച്ചു}.
\end{quote}

\begin{quote}
\emph{kuṭṭikaḷ adhyāpakane anusariccu.}
\end{quote}

\begin{quote}
The children obeyed the teacher.
\end{quote}

Past: \ml{അനുസരിച്ചു}, \emph{anusariccu}.

\ml{അനുസരിച്ച്}, according to, is a form of this verb: following.

\begin{cousinweb}
\textbf{\ml{അനുസരിക്കുക}} (\emph{anusarikkuka}) is built from Sanskrit \emph{anusaraṇa}, \textquotedblleft{}following, going after; conformity\textquotedblright{} (Monier-Williams), with the Malayalam verb ending -\ml{ിക്കുക}. To obey is to follow.
\end{cousinweb}

\begin{grammarlens}[title={What you've built}]
One more word for giving reasons and opinions.
\end{grammarlens}

\subsection*{Your turn}
\begin{itemize}
\raggedright
  \item \emph{Say it:} \emph{anusarikkuka}
  \item \emph{Say it:} \emph{anusarikkuka}, once more
  \item \emph{Say it:} \emph{kuṭṭikaḷ adhyāpakane anusariccu.} — The children obeyed the teacher.
  \item \emph{Recall:} say the Malayalam for a regulation, then the Malayalam for to abide by, then say \emph{anusarikkuka} again
\end{itemize}

\begin{tcolorbox}[breakable,colback=teal!4,colframe=teal!35!black,title={Before you move on}]
What does \ml{അനുസരിക്കുക} mean? (\textquotedblleft{}To obey\textquotedblright{}.) Say it once more.
\end{tcolorbox}
\section[\texorpdfstring{\ml{നിയമപരമായ} (niyamaparamāya)}{നിയമപരമായ (niyamaparamāya)}]{\ml{നിയമപരമായ} (niyamaparamāya) — legal; lawful}

\label{lesson:ML-C417-niyamaparamaya}



\begin{quote}
Before the new one: say the Malayalam for to abide by, then the Malayalam for to obey.
\end{quote}

\subsection*{You'll want to know: \ml{നിയമപരമായ}}
\textbf{\ml{നിയമപരമായ}} — \emph{niyamaparamāya} — \textquotedblleft{}legal; lawful\textquotedblright{}.

\begin{quote}
\ml{ഇത്} \ml{പൂർണ്ണമായും} \ml{നിയമപരമായ} \ml{ഒരു} \ml{കരാറാണ്}.
\end{quote}

\begin{quote}
\emph{itŭ pūrṇṇamāyuṁ niyamaparamāya oru karāṟāṇŭ.}
\end{quote}

\begin{quote}
This is a fully legal contract.
\end{quote}

-\ml{പരമായ}, concerned with, makes adjectives from nouns: \ml{നിയമപരമായ}, to do with the law.

\begin{cousinweb}
\textbf{\ml{നിയമപരമായ}} (\emph{niyamaparamāya}) is built on \ml{നിയമം}, \textquotedblleft{}a law\textquotedblright{}, from Sanskrit \emph{niyama}, \textquotedblleft{}restraint; a fixed law\textquotedblright{} (CDIAL 7250), with -\ml{പരമായ}, from Sanskrit \emph{para}, which at the end of a compound means \textquotedblleft{}devoted to, concerned with\textquotedblright{} (Monier-Williams).
\end{cousinweb}

\begin{grammarlens}[title={What you've built}]
One more word for giving reasons and opinions.
\end{grammarlens}

\subsection*{Your turn}
\begin{itemize}
\raggedright
  \item \emph{Say it:} \emph{niyamaparamāya}
  \item \emph{Say it:} \emph{niyamaparamāya}, once more
  \item \emph{Say it:} \emph{itŭ pūrṇṇamāyuṁ niyamaparamāya oru karāṟāṇŭ.} — This is a fully legal contract.
  \item \emph{Recall:} say the Malayalam for to abide by, then the Malayalam for to obey, then say \emph{niyamaparamāya} again
\end{itemize}

\begin{tcolorbox}[breakable,colback=teal!4,colframe=teal!35!black,title={Before you move on}]
What does \ml{നിയമപരമായ} mean? (\textquotedblleft{}Legal; lawful\textquotedblright{}.) Say it once more.
\end{tcolorbox}
\section[\texorpdfstring{\ml{നിയമവിരുദ്ധമായ} (niyamaviruddhamāya)}{നിയമവിരുദ്ധമായ (niyamaviruddhamāya)}]{\ml{നിയമവിരുദ്ധമായ} (niyamaviruddhamāya) — illegal}

\label{lesson:ML-C417-niyamaviruddhamaya}



\begin{quote}
Before the new one: say the Malayalam for to obey, then the Malayalam for legal.
\end{quote}

\subsection*{You'll want to know: \ml{നിയമവിരുദ്ധമായ}}
\textbf{\ml{നിയമവിരുദ്ധമായ}} — \emph{niyamaviruddhamāya} — \textquotedblleft{}illegal\textquotedblright{}.

\begin{quote}
\ml{അവിടെ} \ml{പാർക്ക്} \ml{ചെയ്യുന്നത്} \ml{നിയമവിരുദ്ധമാണ്}.
\end{quote}

\begin{quote}
\emph{aviṭe pārkkŭ ceyyunnatŭ niyamaviruddhamāṇŭ.}
\end{quote}

\begin{quote}
Parking there is illegal.
\end{quote}

The opposite of \ml{നിയമപരമായ}.

\begin{cousinweb}
\textbf{\ml{നിയമവിരുദ്ധമായ}} (\emph{niyamaviruddhamāya}) is built from Sanskrit \emph{niyama}, \textquotedblleft{}a fixed law\textquotedblright{} (CDIAL 7250), and \emph{viruddha}, \textquotedblleft{}opposed, hindered\textquotedblright{} (CDIAL 11864), with the adjective ending -\ml{ആയ}: against the law.
\end{cousinweb}

\begin{grammarlens}[title={What you've built}]
One more word for giving reasons and opinions.
\end{grammarlens}

\subsection*{Your turn}
\begin{itemize}
\raggedright
  \item \emph{Say it:} \emph{niyamaviruddhamāya}
  \item \emph{Say it:} \emph{niyamaviruddhamāya}, once more
  \item \emph{Say it:} \emph{aviṭe pārkkŭ ceyyunnatŭ niyamaviruddhamāṇŭ.} — Parking there is illegal.
  \item \emph{Recall:} say the Malayalam for to obey, then the Malayalam for legal, then say \emph{niyamaviruddhamāya} again
\end{itemize}

\begin{tcolorbox}[breakable,colback=teal!4,colframe=teal!35!black,title={Before you move on}]
What does \ml{നിയമവിരുദ്ധമായ} mean? (\textquotedblleft{}Illegal\textquotedblright{}.) Say it once more.
\end{tcolorbox}
